\chapter{Protocol for complete-journey generation}
\label{ch:native-protocol}

The protocol first defines the pilot and its fixed-duration control. Sections~\ref{sec:expanded-protocol} and \ref{sec:statistical-perspectives} then specify the expanded replication and distinguish its predeclared endpoints from the later exploratory audit.

\section{Pilot data, partitions and observation units}
The pilot experiment uses native trajectories from the public Monaco SUMO Traffic scenario~\cite{most}. These are simulator observations, not measurements from vehicles on physical roads. The source CNR data motivates the inherited representation and class definitions; the current complete-journey models are fitted on directly extracted SUMO vehicle legs. A bus leg is a completed scheduled service including intermediate stops. Ordinary vehicle legs are bounded by departure, arrival or activity events, rather than selected as the longest convenient moving segment.

The retained vehicle cohort contains 1,051 legs: 925 for training and 126 for development. Each vehicle leg in this cohort has a different parent identity, and parent identities do not cross these partitions. All nevertheless originate from one traffic realization; distinct participants are not independent simulation episodes. Pedestrians are outside the declared five-class vehicle generation task. Table~\ref{tab:native-cohort} reports the actual imbalance instead of implying a balanced dataset.

\begin{table}[htbp]
\centering
\caption{Complete vehicle legs used by both diffusion versions and both GAF readers. The source spelling \texttt{bycicle} is retained in machine-readable labels.}
\label{tab:native-cohort}
\begin{tabular}{lrr}
\toprule
Participant class & Training & Development\\
\midrule
Bus & 81 & 12\\
Motorcycle & 226 & 43\\
Private car & 457 & 50\\
Bicycle & 92 & 14\\
Taxi or ride-hailing vehicle & 69 & 7\\
\midrule
Total & 925 & 126\\
\bottomrule
\end{tabular}
\end{table}

Native quarter-second samples, stopped intervals and long journeys are retained. An event endpoint can occur up to one sampling interval beyond the last observation; no coordinate is invented there. The observed span defines the supervised duration. Map-constructor success is not a training admission criterion: supervised route--motion pairs come directly from the accepted observations. Training-only normalization, duration regression and empirical donors are fitted after selecting training rows. The older frozen test trajectories are not reopened for this experiment's model or parameter selection.

\section{Fixed fits and comparison budget}
The two clock versions each have two independent model fits, with seeds 42 and 43, giving four fits in total. Each receives 10,000 AdamW updates with an effective batch of 64, learning rate $2\times10^{-4}$, weight decay $10^{-4}$ and gradient norm clipping at one. Classes are drawn uniformly, then a training journey is drawn within the chosen class. Length-sorted microbatches obey a 32,768-token budget and a maximum of 32 trajectories. The final exponential moving average, updated with weight 0.001 per step, is the only generation checkpoint; development quality does not select an earlier checkpoint.

Both GAF readers use ResNet18 with ImageNet initialization, seed 42, 50 epochs, batch size 64 and the same AdamW learning rate and weight decay. Cross-entropy class weights are proportional to inverse square-root training frequency. The epoch-50 model is fixed, with no validation checkpoint search, temperature fitting or image augmentation. Accuracy, macro-F1, per-class recall, negative log likelihood, Brier score and 15-bin expected calibration error characterize these components~\cite{calibration}; none is a substitute for generated-trajectory quality.

Each model generates eight candidates per map-ready request using 100 DDIM steps. The empirical comparator selects the eight same-class training routes closest in absolute log length ratio, breaks equal distances by stable identity, and samples one of their whole motion profiles for each candidate. It transfers normalized progress to the common sampled duration and requested route. This preserves training motion structure and supplies a substantive non-neural comparison. The same request and candidate identifiers couple duration draws across all generators. Equal updates or candidate counts do not imply equal computation: valid tokens, memory, elapsed training and sampling time, and teacher calls are reported separately.

\section{A prospective simulator reference}
One new SUMO run provides a reference independent of model fitting. Before simulation, a manifest fixes 100 requests: all 24 public bus services and 19 requests for each of the other four classes. Non-bus requests are selected deterministically from explicit public vehicle plans without activity stops, using source parents assigned to the development partition; source test parents are excluded. They are not class-labelled random origin--destination routes. Concrete vehicle types are drawn deterministically from their source families, preserving actual SUMO vehicle-class permissions: an Uber-labelled passenger vehicle does not acquire taxi-lane access merely through its label.

Endpoints specify a permitted lane and position on the first and last planned edges. The starting position is $\min(20\,\mathrm{m},L_{\mathrm{lane}}/2)$ and the endpoint is the end of the chosen final lane. Public bus stop order and relative schedule semantics are retained. New target departures lie between 18,000 and 18,600 seconds; simulation runs from 14,400 to 23,000 seconds with step 0.25 seconds and seed 330201. Native background demand is retained, while rerouting is disabled only for the new targets. FCD output is restricted to those targets~\cite{sumodocs}. The controlled departures and endpoints make this a new reference experiment, not an exact reconstruction of the source CNR run.

All 100 targets arrived, and their realised edge sequences matched their plans. Predeclared continuity and event checks accepted 97 references: two bus services and one motorcycle were rejected. Independently, the map interface constructed 97 of the 100 planned geometries. These are different success sets: their intersection contains 94 requests before any candidate-level generation or metric-support failure. Coverage therefore always starts at 100, while each quality statistic reports its own paired support. Reference failures are not filled by a new simulation, and unsupported map requests are not replaced by easier ones. One accepted reference contains an emergency stop; it remains in the declared cohort rather than being removed after inspecting quality.

Generation does not consult reference trajectories, arrival times or realized future paths. Reference data are opened for evaluation only after a completed candidate archive is saved. Nevertheless the experiment is a development replication on the same map and public demand templates. A new random seed is not evidence of unseen-line or unseen-city generalization, and shared templates can overlap training geometry.

\section{Physical quality beyond travel distance}
Evaluation computes motion from the returned coordinates and timestamps, independently of the GAF readers. Cumulative planar distance is interpolated onto integer seconds, giving speeds $v_j=d(j+1)-d(j)$ and accelerations $a_j=v_{j+1}-v_j$. Table~\ref{tab:native-endpoints} defines the ten fixed journey summaries. Only complete one-second intervals contribute to these summaries; full observed duration and planar distance retain the final fractional interval.

\begin{table}[htbp]
\centering
\caption{The ten fixed physical summaries. Thresholds define descriptive states rather than asserting vehicle-specific physical limits.}
\label{tab:native-endpoints}
\small
\begin{tabularx}{\textwidth}{@{}lX@{}}
\toprule
Summary & Definition\\
\midrule
Duration; planar distance & Full timestamp span; sum of consecutive planar displacements.\\
Mean; 95th-percentile speed & Mean and empirical 0.95 quantile of $v_j$.\\
Mean absolute acceleration & Mean of $|a_j|$ on the one-second clock.\\
Stopped fraction; longest stop & Fraction and longest run with $v_j<0.3\,\mathrm{m/s}$.\\
Stop episodes per minute & Number of stopped runs divided by observed one-second time, multiplied by 60.\\
Longest steady interval & Longest run with $v_{j+1}\geq0.3\,\mathrm{m/s}$ and $|a_j|<0.5\,\mathrm{m/s^2}$.\\
Four-second speed change & Mean of $|v_{j+4}-v_j|$.\\
\bottomrule
\end{tabularx}
\end{table}

For each class and summary, generated and reference distributions are compared by one-dimensional Wasserstein distance,
\begin{equation}
 W_1(F,G)=\int_{-\infty}^{\infty}|F(x)-G(x)|\,\mathrm{d}x.
 \label{eq:native-wasserstein}
\end{equation}
The unit is the summary's original physical unit. Metrics with different units are not added into an arbitrary overall score. A short journey that cannot support a summary remains in coverage and is marked unsupported for that metric. For $N_c$ jointly supported requests of class $c$, with $K_r$ supported generated candidates in request $r$, candidate mass is $1/(N_cK_r)$ and reference mass is $1/N_c$. Partial candidate failures therefore cannot give a request less weight merely because fewer candidates survived. A five-class macro average is defined only when all five classes have support; available-class summaries are labelled separately.

Native-resolution distributions provide a complementary check. For an interval $\Delta t_i$, speed is $\|X_i-X_{i-1}\|/\Delta t_i$; acceleration divides successive speed differences by $(\Delta t_i+\Delta t_{i-1})/2$. Their Wasserstein distances first weight requests equally, then candidates within a request equally, then observations by their physical interval length. This exposes rapid fluctuations that one-second summaries and the speed GAF can suppress. Generated and reference durations differ, so framewise coordinate error is not treated as the primary synthesis metric.

\section{Diversity, selection and interpretation}
Within a request, candidate pairs are compared using 129 relative-time coordinate positions, 257 relative-time progress values and duration differences. Pair distances are averaged within each request before summarizing across requests. The nearest same-class training progress profile provides a descriptive memorization check, with its associated duration difference. Geometry overlap alone cannot establish memorization when the route was supplied as a condition, and no unvalidated privacy threshold is inferred from the nearest-neighbour statistic.

The main comparison has two panels: all unselected candidates from the empirical and two diffusion generators, followed by Uniform, S, V and SV outputs from the shared pool of each native-clock model. Terminal selection cannot change the quality of its shared pool; its effect concerns the delivered samples. With one selected sample per request, within-pool diversity is no longer estimable, so the selected class distributions, nearest-training comparisons and selection frequencies must be read alongside the original pool diversity.

Both model seeds are reported separately. The single reference episode does not justify confidence intervals that pretend requests are independent simulation replications. Failures, support, coverage and cost accompany quality results. Improved teacher probability alone supports only successful score optimization; useful GAF feedback requires independent motion improvements with any coverage or diversity tradeoff disclosed. These chapters specify that comparison without presuming its outcome or claiming an unperformed downstream augmentation or in-loop guidance result.

\section{A subsequent fixed-duration development control}
After examining the G33 comparison, a bounded follow-up, G34, tests feedback when candidates cannot compete through different total durations. Its protocol is fixed before generating the new pools. For each request and model seed, the predetermined G33 Uniform permutation identifies one candidate slot. Replaying that slot's training-prior draw supplies a single duration $T$ for all eight new candidates. Neither a reference duration nor a teacher score chooses $T$. The same native-clock models, seeds 42 and 43, route conditions, initial-noise keys, $K=8$ and 100 DDIM steps are retained. Other candidates can have a different sequence length because their former duration draws have been replaced; equal candidate and denoising-call counts do not imply equal tensor computation.

The original Uniform slot provides an exact consistency control: its timestamp--coordinate array must remain bitwise identical to the corresponding G33 output. The completed run verifies this for all 97 map-ready requests in each seed, with zero maximum absolute difference. It also verifies that every candidate and each of the four selected outputs uses that request's common duration. No diffusion model, CNN or duration prior is refitted, and no scoring weight or guidance step is searched. Uniform, S, V and SV operate on the same new pool using the unchanged terminal-selection rule.

The follow-up comprises two pools of 776 candidates and eight selected-set evaluations. All retain the original 100-request denominator, 97 generated requests and 94 paired references. The same physical metrics, support rules and diversity diagnostics apply; constant duration across strategies is a control, not a feedback improvement. Because the reference was already analyzed in G33 and motivated this question, G34 is a development mechanism comparison, not a new blind replication. It is reported alongside the original experiment rather than replacing it. Both comparisons concern terminal feedback; neither supplies an in-loop guidance or new downstream-augmentation result.

\section{Expanded-model replication and fixed candidate controls}
\label{sec:expanded-protocol}

The expanded classification audit first evaluates pilot and expanded readers on exactly the same 7,230 development legs. Saved expanded-model predictions are verified against their checkpoint and input hashes; old readers only perform inference. Training parents of both reader versions are disjoint from these development parents. Some old development identities reappear as new stochastic realizations, so this is not a wholly unseen source-template population. Reported parent-bootstrap intervals condition on the existing public domain and fixed models, with 2,000 paired resamples. Leg-weighted results are accompanied by equal-parent sensitivity results.

The generation replication freezes 100 requests per scene before simulator output is available: 24 public bus services and 19 held-out source plans for each other class. The same underlying plans appear in both scenes, with distinct target identities, independent simulator seeds and departure windows of 18,000--18,600 and 32,400--33,000 seconds. The scenes therefore vary traffic context and realization together; they are not an isolated intervention on clock time. Non-bus requests use complete single-leg plans without internal activity parking. The broader training activity-chain support must not be mistaken for a full activity-chain generation evaluation.

Two pilot and two expanded native-time models generate eight candidates per supported request using 100 DDIM steps. One duration is sampled from the corresponding training-side prior and held fixed across the eight candidates. Each model era retains its own prior; the old/new generator comparison consequently tests the complete factorized system. A reference future path, duration or class score never enters sampling. Model seeds, selection rules and experiment limits are locked before outcomes.

For each expanded-generator pool, both old and expanded teacher pairs apply Uniform, spatial S, velocity V and equal-weight SV selection to the exact same trajectories. This is the primary teacher contrast: it holds the generator, candidate budget and duration fixed. Pilot-generator pools additionally receive pilot-teacher selection as a baseline. The unchanged Uniform selections are checked array by array across teacher pairs. Eight full-pool and 48 selected-set quality evaluations result, with no success-based replacement of requests.

\tab{Information and controls in the expanded replication.}{g36-controls}{\begin{tabularx}{\linewidth}{lX}\toprule
Quantity & Fixed or varied \\
\midrule
Reference & Two new SUMO episodes; shared public plans and map \\
Generators & Pilot/expanded native models; seeds 42 and 43 \\
Candidates & Eight per route/class; 100 DDIM steps \\
Duration & Fixed within each pool; sampled from that generator's prior \\
Teacher contrast & Old/new readers applied to the same expanded pool \\
Policies & Uniform, spatial, velocity, equal-weight dual \\
Quality & Independent coordinates; no selecting CNN as ground truth \\
Coverage & All 100 requests retained in each scene denominator \\
\bottomrule\end{tabularx}}

\section{Distributional fidelity and request-level agreement}
\label{sec:statistical-perspectives}

The primary $W_1$ analysis weights classes and requests equally. Candidates divide their request's mass; native intervals are weighted by duration.

For scalar movement summaries, a second, already recorded quantity compares corresponding requests:
\begin{equation}
 \operatorname{MAE}_c(f)=\frac{1}{|Q_c|}\sum_{q\in Q_c}
 \frac{1}{K_q}\sum_{k=1}^{K_q}|f(\widehat{\mathcal T}_{qk})-f(\mathcal T_q)|.
\end{equation}
Here $K_q=1$ after selection. MAE preserves request identity; marginal $W_1$ measures population agreement. One stochastic reference is not a route's unique permissible future, so neither measure alone establishes generative quality.

The post-hoc audit retains every endpoint, averaging five-class macro errors equally over four scene/seed cells before taking percentage changes. Two generator seeds are technical replications, not additional traffic scenes. Improving-cell counts are descriptive. Duration, route length and predominantly fixed mean speed remain controls.
